% Article: Curso Defectologia Vigostky
\documentclass{../../common-shared/cls/Article_Class_RevInfor}

% Search path for images
\graphicspath{{../../common-shared/img/}}

\titulo{Síntese das discussões do curso: A Defectologia em Vigotski: Uma Perspectiva Histórico - Crítica}
\autor{Wagner França Marques}
\local{Brasil}
\data{\today}

% Banner de cabeçalho da Revista InFor, colado ao topo e às laterais do papel
% (full-bleed, margem zero). Impresso automaticamente acima do título.
% Sem esta linha, nenhum banner é desenhado (é opcional por artigo).
\cabecalhoRevInfor{cabecalho-img-revistainfor.png}

\begin{document}

% Pra não precisar ficar escrevento Vigotski com essa variável
% a gente só precisa escrever \Vig{}
\newcommand{\Vig}{Vigotski}



% ---------------------------------------------------------------------------
%  Resumos
%  em português e em inglês, 
%  com as respectivas palavras-chave (parte PRÉ-TEXTUAL: vêm antes de \textual)
% ---------------------------------------------------------------------------
% O ambiente resumo do abnTeX2 imprime o cabeçalho "RESUMO" e o texto do resumo.
% As palavras-chave entram como um parágrafo próprio, com o rótulo em negrito.
% (Conteúdo abaixo é de preenchimento — troque pelo resumo real do artigo.)
\begin{resumo}
Este artigo é o trabalho final do curso A Defectologia de \Vig{} e suas Influências na Educação Especial da CDeP3 - Unesp. O objetivo é elaborar uma síntese de todos os conteúdos e reflexões vivenciados durante o curso.
Assim, este texto vai direto às questões mais relevantes e novas para mim enquanto estudante, evitando obviedades sobre \Vig{} a partir do ponto de vista do aluno do curso, que é o autor deste trabalho.
Este texto parte da apresentação de todos os temas e assuntos relevantes e, em seguida, os retoma com as reflexões e conclusões próprias deste estudante sobre os conteúdos do curso e as principais referências nele utilizadas.
Para sintetizar as ideias principais: foi possível imaginar o cotidiano de \Vig{} atuando como dirigente educacional do governo, sentindo as pressões sociais e políticas de uma época comprometida com os propósitos da Revolução Russa. Foi nesse cenário que \Vig{} produziu suas teorias.
Essas teorias rompem com os antigos paradigmas — o determinismo biológico e o modelo médico da deficiência — em favor de uma abordagem mais social e cultural, denominada Teoria Histórico-Cultural. Nesse contexto, a Defectologia é pensada a partir de um conceito de compensação, no qual cabe ao social a responsabilidade de prover as condições necessárias ao desenvolvimento de todos os seres humanos, inclusive daqueles com deficiência.
Como considerações finais: sou grato pela oportunidade de participar deste curso e aprendi muito sobre a vida e a obra de \Vig{}, bem como sobre a Defectologia como uma abordagem social e cultural da deficiência.
\noindent\textbf{Palavras-chave}: Vigotski. Defectologia. Educação Especial. Desenvolvimento.
\end{resumo}

% Versão em INGLÊS: título traduzido (\tituloestrangeiroRevInfor) seguido do
% ABSTRACT e das Keywords. O ambiente resumo[Abstract] imprime o cabeçalho
% "ABSTRACT"; otherlanguage* garante a hifenização correta em inglês.
\tituloestrangeiroRevInfor{Defectology in Vygotsky: A Historical-Critical Perspective}

\begin{resumo}[Abstract]
\begin{otherlanguage*}{english}
This paper is a final work of the course A Defectology of Vygotsky and its Influences on Special Education of CDeP3 - Unesp. The objective is elaborate a synthesis of all the contents and reflections experienced during the course.
So, this text goes straight to the most relevant and new issues for me as a student, avoiding obviousness about \Vig{} at the point of view of the course student that is this author.
This text starts to present all the relevant themes and subjects and so retake this with this student's own reflections and conclusions about the course contents, and the main references used in the course.
To sintetize the main ideas: Was possible to imagine the \Vig{}'s cotidian life working as a educational government director, feeling the social and political pressures of the time committed with russian revolution purposes. In this scenario \Vig{} produced this theories.
These theories break the old paradigms, the biological determinism and the medical model of disability, in favor of a more social and cultural approach, called Historical-Cultural Theory. In this context the Defectology is thinked where a compensation concept where the social have its responsibility to provide the necessary conditions for the development of all human beings, including those with disabilities.
As the final considerations: I have thankfull for the opportunity to participate in this course, and I have learned a lot about \Vig{}'s life and works, and the Defectology as a social and cultural approach to disability.
\noindent\textbf{Keywords}: Vygotsky. Defectology. Special education. Development.
\end{otherlanguage*}
\end{resumo}



% A partir deste ponto iniciam-se os elementos textuais do artigo, que são obrigatórios. O
% comando \textual imprime o cabeçalho "1. ELEMENTOS TEXTUAIS"
\textual

\section{Introdução}
Este artigo é o trabalho final do curso A Defectologia de Vigotski e
suas Influências na Educação Especial da CDeP3 - Unesp. O objetivo é
atender ao solicitado como trabalho final que é apresentar uma síntese
do conteúdo das relfexões vividas durante o curso.

O curso foi crucial para desvandar o pano de fundo sócio, político e
cultural em que \Vig{} estava comprometido. Esse contexto de revolução
russa requer compreender o pensamento de Marx: O materialistmo
histórico.

\Vig{}, no dia a dia de direção de órgãos públicos de educação sentia
no dia a dia as dores e angústicas de lidar com a distância entre visualizar objetivos
educacionais e alcança-los.

A angústica era maior, provavelmente, a medida que os objetivos
educationais consideravam o protagonismo da escola para alcance de
objetivos não só educacionais, mas também dos objetivos sociais
revolucionários de sua época.

Entre os objetivos educacionais e sociais desejados, se enfrentava
inicialmente o estilo tradicional de ensino. Na linguagem mais popular
pedagógica, caracterizado pela educação bancária, onde o protagonismo
não era do aluno mas do professor ou do conteúdo e também com forte
influência religiosa.

\Vig{} sentiu este cenário, o qual, imagino, tornou-se no combustível
para sua produção teórica que me parece ser uma resposta a esse
cenário de enfrentamento diário em busca de objetivos educacionais e
sociais que pretendia alcançar.

Tal contexto fomentou portanto uma produção teórica mas vivida e
sentida no dia a dia resultando a problematização e apresentação de
conceitos de sua teorica Histórico-Cultural como por exemplo:
Compensação, Supercompensação, rompimento com determinismo biológico
para compreensção da deficiências, identificação do que se entendeu
por defeito primário e secundário dentre outros.

Todo esse assunto apresentado nesta introdução será reapresentado e
problematizado nas demais sessão deste artigo.

As considerações finais deste artigo e também da experiência de
participação no curso é de: agradecimento pela oportunidade;
reconhecimento da qualidade de alto nível do curso, com professores
com histórica de vida envolvida com o tema e quando não diretamente,
professores com vivência no mundo real do chão da escola e que dominam
o assunto; por fim, o curso proveu uma base para que eu pudesse
compreender bastante sobre o materialismo histórico que por sí só já é
bastante complexo.


\section{Produção intelectual de \Vig{} no contexto Materialista Histórico}

Acredito que vou tercer algumas considerações sobre o materialismo
histórico que antes do curso eram impensáveis pra mim. A infra
estrutura, ou seja, as relações econômicas e sociais, definem o mode
de pensar dos indivíduos. Então, o modo de pensar pode ser visto na
super estrutura que são as expressões culturais, institucionais e das
relações sociais em geral.

Por isso, acredito que \Vig{} era menos ideológico e mais prático. Ou
seja, imagino que ele realmente implementava o que ele pensava pra
tentar cimentar seus conceitos na infra estrutura.

Sendo assim, minha percepção foi de que a produção intelectual de
\Vig{} não ocorreu no campo das ideias, mas como resultado da busca
cotidiana na direção de departamentos governamentais da educação e
combates cotidianos às estratégias de ensino tradicionais numa
tentativa de alcançar os objetivos educacionais e consequentemente
sociais.

Sendo assim a impressão é que sua produção não resulta em abstrações
filosóficas ou apenas no campo das ideias mas a partir das
dificuldades do seu cotidiano, ou seja, materialista.

\section{Os objetivos gerais da produção intelectual de \Vig{}}

Para \Vig{} a escola tem o potencial de protagonizar as mudanças
sociais que a revolução pretendia.

Diante disso, é local com maior potencial para o desenvolvimento
humano sendo este portanto o objetivo primordial um conceito
importante pra concepção histórico-cultural.

A escola é, por excelência, o local de maior potencial, justamente
pelas oportunidades mais ricas e intensas de socialização.

Então o processo mediado de socialização é que humaniza e potencializa
o aprendizado.



\section{}


\section{}

\section{}

\lipsum[4-5]

\section{Conclusão}
\lipsum[6]

\postextual
% Shared bibliography (common to all articles) + this article's own references
\bibliography{../../common-shared/bib/references,references}

\end{document}

%%% Local Variables:
%%% mode: LaTeX
%%% TeX-master: t
%%% End:
